\scanpage{19}
\uline{Corollaire} $X'$ normal \inserted{[\unsure{géom\textsuperscript{t}|géom.} \unsure{unibranche} suffit] connexe} de dim $\geq n$,

\noindent\begin{tabular}{@{}|p{.89\linewidth}@{}}
$Y'=X'\cap H_1\cap\cdots\cap H_{\unsure{k|h}}$, où $k\leq n-1$. Alors
\deleted{\unsure{$\pi_1$|$\pi_0$}} $Y'$ est connexe, et $\pi_1(Y')\mathrel{\text{---}}\pi_1(X')$
est surjectif
\end{tabular}

En effet, la connexité : \unsure{déjà|étant} \unsure{pr.|fait|acquise}.
Pour la surjectivité des $\pi_1$, il faut
prouver que si $X'_1$ est un \unsure{rev.|r.v.} étale
\uline{connexe} de $X'$, alors $Y'_1$ est connexe.
Or $X'_1$ \unsure{sera|est} \uline{irréductible}, \deleted{donc} de dim $\geq n$,
d'où le résultat \unsure{\ldots}.

\deleted{\uline{\unsure{Application|Question}}} \uline{\unsure{Conséquence|Conséquences}} Pour vérifier
l'\unsure{engendrement|engendremt|ingrédient} fini des $\pi_1(X')$
dans le cas d'un anneau local $A$,
on est ramené [par \unsure{dévissage|devt}] à le
vérifier pour les anneaux \uline{normaux}
de dim 2. Conjecture : cela \unsure{marche}
\deleted{\unsure{si le corps}} pour les anneaux locaux
\uline{complets}, à corps résiduel \unsure{alg\textsuperscript{t}|alg\textsuperscript{mt}} clos.
% archive sheet number: 18
